% TOP_PROBLEM: {"problemId": "problem.igusa-p-adic-monodromy-conjecture", "canonicalTitle": "Igusa’s Standard p-Adic Monodromy Conjecture", "releaseRank": 367, "primaryDomain": "number_theory_arithmetic_geometry", "primaryDomainLabel": "Number theory & arithmetic geometry", "displayStatus": "open_with_solved_subcases", "releaseStatus": "open_with_solved_subcases"}
% !TeX program = lualatex
% ATTEMPT_STATUS: unresolved
\documentclass[11pt]{article}
\usepackage[margin=25mm]{geometry}
\usepackage{fontspec}
\setmainfont{Latin Modern Roman}
\usepackage{amsmath,amssymb,mathtools}
\usepackage{unicode-math}
\setmathfont{Latin Modern Math}
\usepackage{xurl,hyperref}
\hypersetup{colorlinks=true,urlcolor=blue}
\setcounter{secnumdepth}{0}
\setlength{\parindent}{0pt}
\setlength{\parskip}{5pt}
\setlength{\emergencystretch}{3em}
\begin{document}
\section*{\texorpdfstring{Igusa’s Standard p-Adic Monodromy Conjecture}{Igusa’s Standard p-Adic Monodromy Conjecture}}\label{367}
\subsection{Definitions and mathematical statement}
For nonconstant $f\in{\mathbb{Q}}[x_1,\ldots,x_n]$, prime $p$, and a locally constant compactly supported test function $\phi:{\mathbb{Q}}_p^n\to{\mathbb{C}}$, define for $\Re s>0$
\[
 Z_p(f,\phi;s)=\int_{{\mathbb{Q}}_p^n}|f(x)|_p^s\phi(x){\,\mathrm{d}} x,
 \qquad\operatorname{vol}({\mathbb{Z}}_p^n)=1.
\]
Local complex monodromy at $a\in\{f=0\}$ is the automorphism of the cohomology of the local Milnor fiber induced by moving a small nonzero value once around zero. The conjecture says that outside a finite set of primes depending only on $f$, every actual pole $s_0$ of the meromorphic continuation gives an eigenvalue $e^{2\pi i\Re s_0}$ of such monodromy at some $a$. The excluded-prime set must work for all $\phi$. Candidate poles of a resolution formula that cancel are not poles, and an origin-only or Bernstein--Sato assertion is a different formulation.
\par\smallskip\textit{Formulation and concept references:} \href{https://arxiv.org/pdf/2403.03343}{[S1]}.

\subsection{Short English statement}
Do the actual poles of every relevant $p$-adic zeta integral predict local complex monodromy eigenvalues, outside finitely many primes depending only on the polynomial?
\subsection{Research attempt: arbitrary test functions for monomials and the cancellation barrier}
\textbf{Status.} The standard monodromy implication for arbitrary multivariate polynomials is \textbf{UNRESOLVED} by this attempt. We prove it for monomials in any number of variables and for every one-variable polynomial, with all compactly supported locally constant test functions and every prime. We also give an exact cancellation of an exceptional-divisor candidate pole and a concrete reason that monodromy cannot be checked only at the origin.

\textbf{Source audit.} Veys [S1, Conjecture 2.12(1)] states exactly the standard version: one exceptional-prime set depending on $f$, all test functions, and monodromy at some point of the complex zero locus. Part (2) is the stronger Bernstein--Sato statement; the nearby discussion also distinguishes a local version. These are not interchangeable targets. The elementary subcases below are known-type calculations, not new general monodromy results. Our source search did not verify a general theorem removing all restrictions on the polynomial, so no such theorem is invoked.

\subsubsection{1. The basic one-coordinate integral, including arbitrary-radius balls}
Normalize Haar measure by $\operatorname{vol}(\mathbb Z_p)=1$. For an integer $a\ge1$ and any integer $e$, partition $p^e\mathbb Z_p$ by valuation $k\ge e$. The annulus $v_p(x)=k$ has measure $(1-p^{-1})p^{-k}$, hence for $\Re s>0$,
\[
 I_{a,e}(s)=\int_{p^e\mathbb Z_p}|x|_p^{as}\,dx
 =(1-p^{-1})\sum_{k\ge e}p^{-k(as+1)}
 =\frac{(1-p^{-1})p^{-e(as+1)}}{1-p^{-(as+1)}}.         \tag{1}
\]
The omitted point $x=0$ has measure zero. The expression on the right supplies a meromorphic continuation. Its poles satisfy
\[
 \Re s=-1/a,
\]
with imaginary parts differing by integer multiples of $2\pi/(a\log p)$. The exponential numerator is everywhere nonzero, so these are genuine simple poles for this integral.

If a ball $b+p^e\mathbb Z_p$ does not contain zero, the ultrametric inequality makes $|x|_p=|b|_p$ throughout it. Its integral is
$p^{-e}|b|_p^{as}$, an entire exponential in $s$. If it contains zero, it is exactly $p^e\mathbb Z_p$ and (1) applies. This dichotomy will handle balls centered away from singular points as well as balls around zero.

\subsubsection{2. All test functions for a multivariate monomial}
Take
\[
 f(x)=c\prod_{i=1}^r x_i^{a_i},\qquad
 c\in\mathbb Q^*,\quad a_i\ge1,\quad r\le n.
\]
Let $\phi:\mathbb Q_p^n\to\mathbb C$ be locally constant with compact support. It is a finite linear combination of characteristic functions of pairwise disjoint product balls of a common sufficiently small radius. To justify finiteness, cover the compact support by finitely many balls on which $\phi$ is constant and refine them to a common $p$-adic radius; a compact set meets only finitely many such balls in a bounded ambient ball. Locally constant compact support is itself a union of those balls, so no uncontrolled boundary term remains.

On each product ball, the integral factors by Fubini into the one-coordinate integrals in (1), constant-norm factors for balls not containing zero, and plain volumes for coordinates absent from $f$. Multiplication by $|c|_p^s$ is entire and never zero. Therefore every pole of the full finite sum has real part among
\[
 -1/a_1,\ldots,-1/a_r.                                \tag{2}
\]
Some poles may cancel between summands, especially because $\phi$ may be complex-valued. Cancellation can remove members of the list, but cannot create a pole with a new real part. Repeated exponents can give higher-order poles, without changing (2).

This proves a statement uniform over all $\phi$ for this fixed polynomial and for every prime. The decomposition of $\phi$ may depend on $p$ and on $\phi$; the possible real parts (2) do not. Thus no exceptional-prime set is needed in this monomial subcase. Rational denominators in $c$ only change an entire factor.

\subsubsection{3. The predicted monodromy eigenvalues really occur}
Fix an index $i\le r$, and choose a complex point $a$ with $x_i=0$ and every other coordinate appearing in $f$ nonzero. Near that point,
$f=x_i^{a_i}u(x)$ with $u$ holomorphic and nowhere zero. On a sufficiently small simply connected neighborhood the unit has a holomorphic $a_i$-th root. The coordinate change
$w_i=x_i u(x)^{1/a_i}$, with the other coordinates retained, is locally invertible by the inverse function theorem. Thus the local function is analytically equivalent to $w_i^{a_i}$.

A nearby Milnor fiber has $a_i$ contractible components, one for each root of a small nonzero value. Circling that value around zero permutes the components cyclically. On ordinary $H^0$ the monodromy is the cyclic permutation representation, whose eigenvalues are all $a_i$-th roots of unity. In particular
\[
 \exp(-2\pi i/a_i)
\]
is an eigenvalue. These are exactly the values predicted by (2). The convention includes $H^0$ of the Milnor fiber; when $a_i=1$, eigenvalue one is still present. A reduced-cohomology-only convention would need separate treatment and is not the one used in the catalogue's standard statement.

Combining this local calculation with Section 2 proves the selected implication for every monomial and every test function. The point at which the eigenvalue is seen is allowed to depend on the pole, as the statement requires. We did not insist that all relevant eigenvalues occur at a single chosen singular point.

\subsubsection{4. Why an origin-only replacement is false even for a monomial}
Consider $f(x,y)=x^2y^3$ and $\phi=\mathbf1_{\mathbb Z_p^2}$. The exact zeta function is
\[
 Z_p(f,\phi;s)=
 \frac{(1-p^{-1})^2}{(1-p^{-(2s+1)})(1-p^{-(3s+1)})}.  \tag{3}
\]
Its real poles are $-1/2$ and $-1/3$, predicting $-1$ and $e^{-2\pi i/3}$. They occur along the smooth parts of the coordinate axes as in Section 3.

At the origin the local monodromy has only eigenvalue one. This can be seen explicitly, without guessing from the number of components. For a nonzero fiber value $t$, set $z=xy$. Then
\[
 x=t^{-1}z^3,\qquad y=t z^{-2}.                        \tag{4}
\]
These formulas give inverse parametrizations of $x^2y^3=t$ by $z\in\mathbb C^*$. Its intersection with a sufficiently small Euclidean ball about the origin is described by
\[
 |t|^{-2}|z|^6+|t|^2|z|^{-4}<\epsilon^2.
\]
For sufficiently small $|t|>0$, the allowed radii form a nonempty annulus. The formula depends only on $|t|$, so as $t$ travels around a circle, keeping $z$ fixed in (4) gives a trivialization returning identically after one turn. Thus the monodromy acts as the identity on both $H^0$ and $H^1$ of this annulus; other cohomology vanishes.

The calculation does not contradict the catalogue formulation, because that formulation permits monodromy at some point of $\{f=0\}$. It instead gives a complete diagnostic against silently strengthening the target to the origin alone. It also distinguishes local topology at a crossing from that at a general point of one component.

\subsubsection{5. Every one-variable polynomial, every prime and every test function}
Let $f\in\mathbb Q[x]$ be nonconstant. For fixed $p$, it has finitely many roots $a\in\mathbb Q_p$. Near each such root,
\[
 f(x)=(x-a)^{m_a}g_a(x),\qquad g_a(a)\ne0.
\]
On a sufficiently small ball, $|g_a(x)|_p$ is constant by continuity and the ultrametric inequality. Shrink the ball further so that $\phi$ is constant there. Its contribution is a constant times an entire exponential times $I_{m_a,e}(s)$, and hence has possible real pole $-1/m_a$.

Remove these finitely many root neighborhoods from the compact support. On the remaining compact set, $f$ has no zeros. The positive function $|f|_p$ is locally constant there and takes only finitely many values, again by compactness. The remaining integral is a finite sum of exponentials in $s$ and is entire. Thus every actual pole has real part $-1/m_a$ for some $p$-adic root.

Such a root is algebraic over $\mathbb Q$, since it is a root of the fixed rational polynomial. Its multiplicity is preserved under an embedding of the number field $\mathbb Q(a)$ into $\mathbb C$. At the resulting complex root, the local function is analytically equivalent to $w^{m_a}$ and has all $m_a$-th roots of unity as monodromy eigenvalues. The implication follows exactly as before.

This argument works for every prime, including primes dividing coefficient denominators or discriminants; the neighborhoods may become smaller but the logic is unchanged. It proves the standard conjecture in one variable with an empty exceptional-prime set, uniformly for all test functions. It does not use a complex embedding of $\mathbb Q_p$ as a topological field; it only embeds the finite number field generated by an algebraic root.

\subsubsection{6. Resolution denominators are candidates, not conclusions}
In a normal-crossings chart of an embedded resolution one has expressions of the form
\[
 f\circ h=u\prod_i y_i^{N_i},\qquad
 |\operatorname{Jac}h|=|v|\prod_i|y_i|^{\nu_i-1},
\]
with nonvanishing local units $u,v$. Summing valuation annuli produces denominator factors
\[
 1-p^{-(N_i s+\nu_i)}.                                \tag{5}
\]
Their real candidate poles are $-\nu_i/N_i$. This computation generalizes (1), but a sum over charts and strata can cancel such factors. It is not legitimate to conclude that every exceptional divisor contributes an actual pole or a required new monodromy eigenvalue.

A complete elementary cancellation is visible by resolving the already smooth function $f(x,y)=x$ through a blowup at the origin. Partition $\mathbb Z_p^2$, ignoring measure-zero axes, into the regions $v_p(x)\le v_p(y)$ and $v_p(y)<v_p(x)$. In the first put $y=xt$, with $t\in\mathbb Z_p$; in the second put $x=yt$, with $t\in p\mathbb Z_p$. The Jacobians have absolute values $|x|_p$ and $|y|_p$, respectively. The two contributions are
\[
 A(s)=\frac{1-p^{-1}}{1-p^{-(s+2)}},\qquad
 B(s)=\frac{(1-p^{-1})^2p^{-(s+1)}}
             {(1-p^{-(s+2)})(1-p^{-(s+1)})}.
\]
But their sum is
\[
 A(s)+B(s)=\frac{1-p^{-1}}{1-p^{-(s+1)}},              \tag{6}
\]
as is also immediate by integrating $|x|_p^s$ directly. The exceptional candidate $s=-2$ has canceled exactly. The common numerator contains $1-p^{-(s+2)}$ after addition.

The blowup was unnecessary geometrically, but useful adversarially: a proof must not depend on every denominator appearing in a chosen resolution being genuine. Different resolutions can introduce extra exceptional divisors while leaving the integral unchanged.

\subsubsection{7. The missing topological bridge after resolution}
For a monomial at a general point of a component, the integer controlling the pole and the cyclic permutation of Milnor-fiber components are visibly the same exponent. On a resolution of an arbitrary singularity, the ratio in (5) involves both vanishing order $N_i$ and Jacobian discrepancy $\nu_i$. The desired eigenvalue has order dividing $N_i/\gcd(N_i,\nu_i)$, but the existence of a divisor with that numerical ratio is not itself a nonvanishing cohomology statement.

Monodromy zeta formulas combine multiplicities with Euler characteristics of resolution strata. Those Euler characteristics may vanish or cancel; moreover an actual zeta-integral pole can be represented by contributions from several divisors. A general proof must relate the noncancellation of the arithmetic integral to noncancellation in the relevant local monodromy data. Our local-coordinate calculation does not establish that relationship.

Nor can the exceptional-prime quantifier be repaired by proving a theorem separately for each $\phi$ with a different excluded set. The catalogue requires one finite set for the fixed polynomial before $\phi$ is chosen. In the subcases proved here the empty set works, which resolves this issue completely there. In a general resolution argument, any good-reduction or comparison theorem would have to be uniform in the test function.

\subsubsection{8. Verification and final conclusion}
The saved diagnostic checks the exact rational identity (6) with $T=p^{-s}$, the pole exponents in (3), and residue-count approximations to the one-variable annulus sum at fixed positive integer $s$. These checks use exact integer or rational arithmetic. They do not infer a pole from a finite list of congruence counts, nor infer a monodromy spectrum from a floating-point root calculation.

The complete proved subcases are all monomials in arbitrary dimension and all one-variable rational polynomials, for every prime and every permitted test function. The origin-only variant fails in the displayed two-variable example, while the selected some-point formulation is verified. The first remaining general step is the arithmetic-to-topological noncancellation implication after resolution, with uniform exceptional primes. The standard multivariate conjecture is \textbf{UNRESOLVED} by this attempt.

\subsection{Sources}
\begin{itemize}
\item[S1]Willem Veys, Introduction to the monodromy conjecture, arXiv:2403.03343v1, Conjecture 2.12(1). Denef, Bourbaki exposé 741, supplies historical corroboration with a separately mentioned test-function extension.
\url{https://arxiv.org/pdf/2403.03343}.
\textit{Formulation source.}
\end{itemize}
\end{document}
